\documentclass[10pt,a4paper]{amsart}
\usepackage[margin=19mm]{geometry}
\usepackage{amsmath,amssymb,mathtools}
\usepackage[scr=rsfs,cal=euler]{mathalfa}
\usepackage{xfrac}
\usepackage[unicode,hidelinks,bookmarks=false]{hyperref}
\newcommand{\ie}{i.e.\ }
\newcommand{\cf}{cf.\ }
\newcommand{\resp}{respectively\ }
\newcommand{\Subsection}{\subsection}
\providecommand{\nobreakdash}{\nobreak\hbox{-}}
\setlength{\emergencystretch}{2em}
\multlinegap0pt
\usepackage{bm}
\usepackage{bbm}
\usepackage{dsfont}
\usepackage{xspace}
\usepackage{cancel}
\usepackage{mathtools}
\usepackage{cleveref}
\usepackage{amsthm}
\makeatletter
\@ifundefined{theorem}{\newtheorem{theorem}{Theorem}}{}
\@ifundefined{lem}{\newtheorem{lem}[theorem]{Lemma}}{}
\makeatother
\DeclareMathOperator{\BS}{BS}
\DeclareMathOperator{\Mon}{Mon}
\newcommand{\Z}{\mathbb{Z}}
\title[The Diophantine problem in the Baumslag-Gersten group]{The Diophantine problem in the Baumslag-Gersten group}
\author{Robert D. Gray, Alex Levine}
\date{Source: arXiv:2607.24189v1 (CC BY 4.0)}
\begin{document}
\maketitle

\noindent\emph{Source and licence.} The Diophantine problem in the Baumslag-Gersten group. Version arXiv:2607.24189v1 (CC BY 4.0), distributed under the Creative Commons Attribution 4.0 International licence, as recorded in the arXiv licence metadata for this exact version and shown on the abstract page https://arxiv.org/abs/2607.24189v1. Full rights notice: licenses/arxiv-2607.24189-CC-BY-4.0.txt.\par\smallskip

\noindent\textit{Editorial scope.} Counted unit: the Lem whose source label is \texttt{defable} in the article (source0.tex lines 894--914, source label \texttt{defable}), delivered together with the article's own complete original proof of it (source0.tex lines 916--1033, one \texttt{proof} environment of 118 lines). No mathematics is written, repaired or re-derived: statement and proof are the article's own text line for line. Required context retained uncounted: BS-cents (lines 340--354); item:BS-cent-a (lines 344--353); item:BS-cent-abi (lines 344--353); item:BS-cent-b (lines 344--353); conj-into-ab (lines 653--659); inf-pres (lines 726--737); cents (lines 782--787). Every definition, display and label of those lines is retained unchanged. Omitted from this package and not redistributed: 1--339, 354--652, 660--725, 738--781, 788--893, 915--915, 1034--1198. Nothing inside a retained excerpt is omitted or altered: every retained body is delivered exactly as the article's lines are written, so no omission receipt of any kind accompanies this package. Citations: the retained excerpts carry no citation command at all, so no bibliography entry of the article is transcribed and no reference is redistributed. Reference adaptation: none was needed and none was made; every reference inside the delivered unit resolves inside it, so no unresolved reference or citation is produced. Printed slips kept literal: no printed word, symbol, hypothesis or number of the counted statement or of its proof was altered. Layout and class: the article's own class is \texttt{amsart} with packages including amssymb, amsmath, amsthm, mathrsfs, faktor, tikz, thmtools, mathtools, hyperref, inputenc, enumitem, wasysym, cleveref; the wrapper is the lane's pinned \texttt{amsart} editorial wrapper at 10pt on A4, which restates the theorem-like environments the retained text needs and adds the article's own macro definitions that the retained text needs (\texttt{\textbackslash BS}, \texttt{\textbackslash Mon}, \texttt{\textbackslash Z}), with extra wrapper packages (bm, bbm, dsfont, xspace, cancel, mathtools, cleveref). The delivered numbering is the unit's own. Assets: no retained span contains a figure environment, a graphics-inclusion command, a PDF-page inclusion, a table or a TikZ picture; no image file is copied into this package, the package declares no asset dependency and no image file is redistributed with it. Licence: version arXiv:2607.24189v1 of this article is distributed under the Creative Commons Attribution 4.0 International licence, as recorded in the arXiv licence metadata for this exact version and shown on the abstract page https://arxiv.org/abs/2607.24189v1. A full rights notice is delivered at licenses/arxiv-2607.24189-CC-BY-4.0.txt. Characters: every retained excerpt is pure ASCII, so the delivered engine, XeLaTeX with its default Latin Modern text font, reports no missing character.
\begin{lem}
	\label{BS-cents}
	Let \(p \in \Z\) be such that \(p > 1\). In \(\operatorname{BS}(1, p) =
	\langle a, b \mid bab^{-1} = a^p \rangle\), we have:
	\begin{enumerate}
		\item \(C_{\operatorname{BS}(1, p)}(a) = \langle b^{-k} a b^k \; (k \in \Z) \rangle\);
			\label{item:BS-cent-a}
		\item \(C_{\operatorname{BS}(1, p)}(b^i) = \langle b \rangle\) for all
				\(i \in \Z \setminus \{0\}\);
			\label{item:BS-cent-b}
		\item \(C_{\operatorname{BS}(1, p)}(ab^i) = \langle ab^i \rangle\) for all
				\(i \in \Z \setminus \{0\}\).
			\label{item:BS-cent-abi}
	\end{enumerate}
\end{lem}

\begin{lem}
		\label{conj-into-ab}
		Let \(H_p = \langle a_i \; (i \in \Z) \mid a_{i + 1} a_i a_{i + 1}^{-1} 
		= a_i^p \; (i \in \Z)\rangle\). If \(h \in \langle a_2^i a_1 a_2^{-i} \;
		(i \in \Z) \rangle\) is such that \(h a_0 h^{-1} \in \langle a_0, a_1
		\rangle\), then \(h \in \langle a_1 \rangle\).
\end{lem}

\begin{lem}
	\label{inf-pres}
	Let \(G_p = \langle a, b, t \mid bab^{-1} = a^p, tat^{-1} = b\rangle\), where
	\(p \in \Z\) is such that \(p > 1\). Let
	\[
	 H_p = 
			\langle a_i \; (i \in \Z) \mid a_{i + 1} a_i
	a_{i + 1}^{-1} = a_i^p \; (i \in \Z)\rangle.
	\]
	Then \(G \cong H_p \rtimes_\varphi \Z\), where \(\varphi\) is defined by
	\(\varphi(a_i) = a_{i + 1}\) for all \(i \in \Z\).
\end{lem}

\begin{lem}
		\label{cents}
		Let \(G_p = \langle a, b, t \mid bab^{-1} = a^p, tat^{-1} = b\rangle\),
		where \(p \in \Z\) is such that \(p > 1\).  Then \(C_{G_p}(a) = \langle
		b^{-i} a b^i \; (i  \in \Z_{\geq 0}) \rangle\).
\end{lem}

\begin{lem}
	\label{defable}
	The following sets are definable in the positive existential theory of
	\(G_p = \langle a, b, t \mid bab^{-1} = a^p, tat^{-1} = b\rangle\): 
	\begin{enumerate}
			\item \(\langle b \rangle\);
					\label{item:b-defable}
			\item \(\langle a \rangle\);
					\label{item:a-defable}
			\item \(\Mon \langle b \rangle\);
					\label{item:mon-b-defable}
			\item \(\langle a, b \rangle\);
					\label{item:ab-defable}
			\item \(\{(b^i, b^{ij}) \mid i, j \in \Z\}\);
					\label{item:b-div-defable}
			\item \(\{(b^i, b^{i p^r}) \mid i \in \Z, r \in \Z_{\geq 0}\}\);
					\label{item:b-divp-pre-defable}
			\item \(\{(b^i, b^{\pm i p^r}) \mid i \in \Z, r \in \Z_{\geq 0}\}\).
					\label{item:b-divp-defable}
	\end{enumerate}
\end{lem}

\begin{proof}
 	\eqref{item:b-defable}: We first apply \cref{cents} to show that
	\[
			C_{G_p}(b) = C_{G_p}(tat^{-1}) = t C_{G_p}(a) t^{-1}
			= t \langle b^{-i} a b^i \; (i \in \Z_{\geq 0}) \rangle t^{-1}
					= \langle  (tbt^{-1})^{-i} b (tbt^{-1})^i \; (i \in \Z_{\geq 0}) \rangle.
	\]
	We then consider the system \((Xb = bX) \wedge (XaX^{-1} a = a XaX^{-1})\).
	The first equation insists that a solution \(x\) to \(X\) lies in \(\langle
	(tbt^{-1})^{-i} b (tbt^{-1})^i \; (i \in \Z)\rangle\), and the second
	equation insists (again, using \cref{cents}) that \(xax^{-1} \in C_{G_p}(a)
	=  \langle b^{-i} a b^i \; (i \in \Z_{\geq 0}) \rangle\). Using the
	description of \(G_p\) as \(H_p \rtimes_\varphi \Z\) from
	\cref{inf-pres}
	along with
	\cref{conj-into-ab}, the only elements of \(\langle (tbt^{-1})^{-i} b (tbt^{-1})^i \;
	(i \in \Z)\rangle\) that conjugate \(a\) into \(\langle a, b \rangle\) are
	those in \(\langle b \rangle\). Thus every solution to this system of
	equations lies in \(\langle b \rangle\).  Since conjugates of \(a\) by
	powers of \(b\) commute with \(a\), it follows that \(\langle b \rangle\)
	is precisely the set of solutions.

	\eqref{item:a-defable}: Since \(\langle b \rangle\) is equationally
	definable by \eqref{item:b-defable}, and \(t^{-1} b^i t = a^i\) for all \(i
	\in \Z\), it follows that \(\langle a \rangle\) is the set of solutions to
	\(Y\) in the system \((X \in \langle b \rangle) \wedge (t^{-1} X t = Y)\).

	\eqref{item:mon-b-defable}: Note that from \eqref{item:a-defable} and
	\eqref{item:b-defable}, \(\langle a \rangle\) and \(\langle b \rangle\) are
	equationally definable in \(G_p\). Then consider the system \((X \in
	\langle b \rangle) \wedge (Y \in \langle a \rangle) \wedge (X aX^{-1} =
	Y)\). Noting that \(\langle a, b \rangle\) is a copy of \(\BS(1, p)\), we
	can use the fact that \(b^i a b^{-i} \in \langle a \rangle\) if and only if
	\(i \geq 0\), and thus it follows that the set of solutions to \(X\) in
	this system is \(\Mon\langle b \rangle\).

	\eqref{item:ab-defable}: Using \cref{cents} and \eqref{item:b-defable}, 
	we have that \(C_{G_p}(a) = \langle b^{-i} a b^i \; (i \in \Z_{\geq 0}) \rangle\)
	and \(\langle b \rangle\) are both equationally definable in \(G_p\).
	Recall that \(\langle a, b \rangle\) is a solvable Baumslag-Solitar group,
	and is thus the semidirect product of these subgroups. In particular,
	\(\langle a, b \rangle = \langle b^{-i} a b^i \; (i \in \Z_{\geq 0}) \rangle
	\langle b \rangle\). As a product of equationally definable sets,
	\(\langle a, b \rangle\) is thus equationally definable.

	\eqref{item:b-div-defable}: Using \eqref{item:ab-defable}, the solvable
	Baumslag-Solitar subgroup \(B = \langle a, b \rangle\) is equationally
	definable. It thus suffices to show that \(\{(b^i, b^{ij}) \mid i, j \in
	\Z\}\) is equationally definable in \(B\). We have from \cref{BS-cents}
	\eqref{item:BS-cent-b} that \(C_B(b) = \langle b \rangle\). Thus the
	set of solutions to the system of equations in \(B\): \((Xb = bX) \wedge (Y = bX^2)\) is
	\(\{(b^i, b^{2i + 1}) \mid i \in \Z\}\). It follows that
	\(\{b^i \mid i \in 2 \Z + 1\}\) is equationally definable in \(B\). From
	\cref{BS-cents}\eqref{item:BS-cent-abi}, we have that \(C_B(ab^i) = \langle ab^i
	\rangle\) for all \(i \in 2 \Z + 1\). Thus the set of solutions to
	\((Y \in \{b^i \mid i \in 2\Z + 1\}) \wedge (Z (aY) = (aY)Z)\) is
	\[
		\{(b^i, (ab^i)^j) \mid i, j \in \Z, i \text{ odd}\}.
	\]
	We have that for all \(i, j \in \Z\), \((ab^i)^j = a^{1 + p^i + \cdots +
	p^{(j - 1)i}} b^{ij}\). Thus if we add the equations \(Ua = aU\), \(Vb =
	bV\) and \(UV = Z\), (using the fact that \(B\) is a semidirect product of
	\(C_B(a)\) with \(\langle b \rangle\) from
	\cref{BS-cents}\eqref{item:BS-cent-a}), we have that the set of solutions
	to \((Y, Z, U, V)\) in this system will be
	\[
			\{(b^i, a^{1 + p^i + \cdots + p^{(j - 1)i}} b^{ij},  a^{1 + p^i +
			\cdots + p^{(j - 1)i}}, b^{ij}) \mid i, j \in \Z, i \text{ odd}\}.
	\]
	In particular, the set \(\{(b^i, b^{ij}) \mid i, j \in \Z, i \text{ odd}\}\)
	is equationally definable in \(B\). It remains to `remove' the oddness
	of \(i\). To do so, we take \((Y, V)\) from the previous system, and add
	the equations
	\[
		(\bar{Y} b = b \bar{Y}) \wedge
		(\bar{V} b = b \bar{V}) \wedge
		(b \bar{Y}^2 = Y) \wedge
		(b \bar{V}^2 = V).
	\]
	The first two equations insist all solutions to \(\bar{Y}\) and \(\bar{V}\)
	will lie in \(\langle b \rangle\) and the latter two insist that the exponents of
	the solutions to \((Y, V)\) can be obtained from the exponents of the solutions
	to \((\bar{Y}, \bar{V})\) by doubling and then adding \(1\). In other words,
	the set of solutions to \((\bar{Y}, \bar{V}, Y, V)\) is
	\[
			\{(b^{\bar{i}}, b^{\bar{i} j}, b^{2 \bar{i} + 1},
			b^{(2 \bar i + 1)j}) \mid \bar{i}, j \in \Z\}.
	\]
	We have thus shown the desired set is equationally definable in \(B\), and thus it
	is also equationally definable in \(G_p\), since \(B\) is equationally definable
	in \(G_p\).

	\eqref{item:b-divp-pre-defable}: First note that \(\langle a \rangle\) and
	\(\Mon\langle b \rangle\) are equationally definable by \eqref{item:a-defable}
	and \eqref{item:mon-b-defable}. Then consider the system \((X \in \langle a
	\rangle) \wedge (Y \in \Mon\langle b \rangle) \wedge (Z = YXY^{-1}) \). Since
	\(b^r a^i b^{-r} = a^{ip^r}\) for all \(i \in \Z\) and \(r \in \Z_{\geq 0}\) 
	the set of solutions to \((X, Y, Z)\) in this system is
	\[
			\{(a^i, b^r, a^{ip^r}) \mid i \in \Z, r \in \Z_{\geq 0}\}.
	\]
	We have thus shown that \(\{(a^i, a^{ip^r}) \mid i \in \Z, r \in \Z_{\geq 0}\}\)
	is equationally definable in \(G_p\). It thus follows that if
	\((X, Y)\) are variables taking solutions in this set, and we add new
	equations \((X' = tXt^{-1}) \wedge (Y' = t Y t^{-1})\), then the
	solutions to \((X', Y')\) will be the desired set.

	\eqref{item:b-divp-defable}: Note that the system
	\[
			(X, Y) \in \{(b^i, b^{i p^r}) \mid i \in \Z, r \in \Z_{\geq 0}\}) \wedge
			(Z = Y^{-1})
	\]
	has a solutions set \(\{(b^i, b^{i p^r}, b^{-ip^r}) \mid i \in \Z, r \in \Z_{\geq 0}\}\),
	and so \(\{(b^i, b^{-ip^r}) \mid i \in \Z, r \in \Z_{\geq 0}\}\) is equationally
	definable in \(G_p\). It thus follows that the desired set is definable in the
	positive existential theory of \(G_p\), as a union of two equationally definable
	sets in \(G_p\).
\end{proof}

\end{document}
